% COLM 2026 workshop submission: "AI Measurement Science: Toward Rigorous AI Evaluation".
% Forked from the ACL/EMNLP source; the frozen ACL version remains at ../main.tex.
% Anonymity + line numbers are controlled by the colm2026_conference package option:
%   [submission] = double-blind anonymous + line numbers (default, for review)
%   [preprint]   = named authors, no line numbers
%   [final]      = camera-ready, named authors
% build-colm.sh swaps the option; do not hand-edit it here.
\documentclass{article}
\usepackage[submission]{colm2026_conference}

% colm2026_conference.sty already provides the mandated fonts (tgpagella/mathpazo/
% inconsolata) and loads natbib + fancyhdr. Do NOT load times or override fonts.
\usepackage[T1]{fontenc}
\usepackage[utf8]{inputenc}
\usepackage{microtype}
\usepackage{amsmath}
\usepackage{amssymb}
\usepackage{graphicx}
\usepackage{booktabs}
\usepackage{multirow}
\usepackage{array}
\usepackage{lineno}
\usepackage{hyperref}
\usepackage{url}
\usepackage{listings}
\lstset{
  basicstyle=\ttfamily\footnotesize,
  breaklines=true,
  breakatwhitespace=false,
  columns=flexible,
  keepspaces=true,
  frame=single,
  xleftmargin=4pt,
  xrightmargin=4pt,
}

\definecolor{darkblue}{rgb}{0, 0, 0.5}
\hypersetup{colorlinks=true, citecolor=darkblue, linkcolor=darkblue, urlcolor=darkblue}

\title{Measuring the Measurement: Repeated-Run Reliability of LLM Code Evaluation}

% Author block is shown only in [preprint]/[final]; [submission] anonymizes automatically.
\author{Yongxi Zhou,$^{1}$ Lai Yun Choi,$^{1}$ Jiaxi Wen,$^{1}$ Wenbo Ye,$^{2}$ \& Junwei Yao$^{1}$ \\
  $^{1}$Northeastern University, Massachusetts, USA \\
  $^{2}$University of Southern California, California, USA \\
  \texttt{\{zhou.yongx, choi.lai, wen.jiax, yao.junw\}@northeastern.edu}, \texttt{ywb6113@gmail.com}}

\begin{document}

\ifcolmsubmission
\linenumbers
\fi

\maketitle

\begin{abstract}
Evaluation scores are measurements, and like any measurement they have a reliability that is rarely quantified. We ask whether single-run pass rate---the dominant instrument for large language model (LLM) code evaluation---is a reliable measurement of model capability, or whether it conceals substantial test--retest variation. Treating each problem as a repeated measurement, we introduce a repeated-run protocol with three instruments: run-level pass rate (RLPR), perfect stability rate (PSR, the fraction of tasks solved on \emph{every} retry), and average per-problem variance (AV). On a recency-controlled benchmark of 100 LeetCode-style problems, we evaluate 16 models from five provider families under two prompt templates with five repeated runs each (16{,}000 evaluation instances). Although RLPR and PSR are strongly correlated ($r=0.985$), run-level pass rate overstates retry-free coverage by up to 17.8 percentage points, the gap is largest for mid-performing systems, and it reverses model rankings even among closely matched systems. Prompt effects are model-dependent rather than uniformly beneficial. We argue that reporting the test--retest reliability of a score is a necessary part of valid LLM evaluation, with deterministic single-turn tasks as the tractable i.i.d.\ baseline for harder non-stationary settings.
\end{abstract}

\section{Introduction}
An evaluation score is a measurement, and every measurement instrument has a \emph{reliability}: the degree to which it returns the same value when the same quantity is measured again. For large language models (LLMs), the dominant instrument---single-run pass rate on a held-out benchmark---is almost always reported as a point estimate, as if it were perfectly repeatable. Yet LLM decoding is stochastic, so re-running the same model on the same task is a \emph{repeated measurement} that need not agree with itself. Borrowing the language of psychometrics and metrology, we ask a measurement-science question: what is the \emph{test--retest reliability} of pass-rate evaluation, and how much does a single run mislead us about a model's underlying capability? We study this in the cleanest available setting---deterministic, single-turn programming tasks graded by a fixed verifier---where the only source of disagreement between repeated measurements is the model itself.

Large language models are increasingly used for text-conditioned tasks with deterministic specifications: given a fixed textual input and a fixed verifier, correctness is binary. Program synthesis and code assistance are especially important instances of this broader setting, because they pair natural-language problem statements with executable test suites. Conventional benchmarks typically report single-run performance, or summarize repeated sampling through eventual-success metrics such as pass@$k$ \citep{chen2021codex,austin2021program}. Recent work has shown that repeated sampling can be scaled as a form of inference-time compute \citep{brown2024large}, yet the \emph{stability} of outcomes across such repeated invocations---whether a model consistently succeeds or fails on the same task---has received less systematic attention. In practice, repeated calls under the same task description can yield different outputs and different pass/fail outcomes, even when the downstream workflow expects reproducible behavior \citep{atil2024llm,alvarado2025repetitions}.

This paper studies that repeated-invocation variability directly. Rather than treating repeated sampling only as an opportunity to recover one success among many attempts, we investigate the relationship between correctness and stability under fixed conditions. Our focus is therefore complementary to pass@$k$ \citep{chen2021codex}: whereas pass@$k$ asks whether a model can produce \emph{at least one} correct solution in $k$ attempts---a capability measure---we ask how reliably it succeeds on \emph{every} invocation, a deployment reliability measure. These two questions target different downstream concerns and can yield divergent rankings, particularly for mid-accuracy systems.

The paper is organized around three measurement-validity questions:
\begin{enumerate}
  \item How large is the gap between a model's run-level pass rate and the fraction of tasks it solves reliably on \emph{every} retry---that is, how much measurement error does single-run reporting hide?
  \item Which aggregate summaries best expose the per-task measurement error (test--retest variability) that a single pass-rate number conceals?
  \item As an exploratory analysis: do design choices---prompt template, reasoning-specialized architectures---systematically change the reliability of the measurement?
\end{enumerate}

We focus on deterministic tasks because they represent a broader and practically important class of NLP system behavior: the model is conditioned on text, but the output is consumed by a verifier that expects reproducible correctness. Code-generation systems are a particularly stringent instance of this pattern, alongside tasks such as text-to-SQL \citep{gao2023text}, structured semantic parsing, tool-call argument generation, and formal text transformations. In such settings, the accuracy--stability gap has direct practical consequences: a model with 60\% run-level pass rate but only 47\% perfect stability solves many individual attempts, yet still leaves a substantial fraction of tasks in a retry-required regime. Our aim is therefore not merely to rank models by mean performance, but to characterize how much of their observed accuracy translates into reliable repeated-use behavior. We deliberately restrict attention to deterministic, single-turn tasks because they are the tractable i.i.d.\ case of repeated measurement: each run is an independent draw under a fixed verifier. This is the natural baseline for the harder, non-stationary regimes---interactive, path-dependent, or adaptively optimized evaluation---where measurement reliability is even less understood, and where the analysis tools developed here would need to be extended.

\section{Related Work}
Code-generation evaluation is commonly framed around functional correctness and pass@$k$ under automated testing \citep{chen2021codex,austin2021program}. Foundational benchmarks emphasize whether a model eventually produces at least one correct program under repeated sampling \citep{chen2021codex,austin2021program,hendrycks2021apps}, while more recent efforts address contamination and temporal validity \citep{jain2024livecodebench}, with parallel work extending evaluation to real-world repository-level issues \citep{jimenez2024swebench}. Prompting and inference-time interventions use additional computation to improve performance rather than to measure it: zero-shot chain-of-thought \citep{kojima2022zeroshot} restructures the prompt, while self-consistency and adaptive-consistency \citep{wang2022self,aggarwal2023lets} aggregate multiple samples. In parallel, work on LLM output non-determinism shows that performance varies substantially across repeated runs under identical conditions \citep{atil2024llm,alvarado2025repetitions}, and prompt sensitivity research demonstrates that small phrasing changes can shift accuracy \citep{zhuo2024prosa}.

Perfect Stability Rate (PSR) is mathematically equivalent to pass@$k$ at $k=R$ (the number of repeated runs) with an all-successes threshold; the distinction is interpretive. Pass@$k$ estimates $P(\text{at least one success in }k\text{ draws})$---a \emph{capability} measure---while PSR estimates $P(\text{all }k\text{ draws succeed})$---a \emph{deployment reliability} measure. These quantities diverge most sharply for mid-accuracy systems. The contribution lies in treating PSR and Average Variance (AV) as first-class evaluation targets for deterministic code tasks---both defined formally in Section~\ref{sec:metrics}---and using them to characterize the accuracy--stability gap directly.

Recent work emphasizes that code evaluation is sensitive to benchmark construction and reporting choices \citep{liu2024your,du2024evaluating}. Our framing shifts the emphasis from the \emph{construction} validity of a benchmark to the \emph{reliability} of the measurement it produces: we treat an evaluation as an instrument whose test--retest behavior can be quantified, in the spirit of reliability theory in psychometrics and metrology, rather than reported as a single repeatable number. Work on confidence calibration \citep{guo2017calibration} and performance--efficiency trade-offs \citep{zhang2026performance} reflects a broader recognition that average accuracy alone is an incomplete measure of model reliability. This perspective is deliberately narrower than end-to-end agent evaluation: it isolates repeated-run behavior for single-turn coding, while broader orchestration, tool use, and runtime dependencies require separate analysis.

\section{Experimental Setup}

\subsection{Problem Dataset}
We evaluate on a curated dataset of deterministic programming problems, denoted $\mathcal{D}=\{(x_i, T_i)\}_{i=1}^{N}$, where $x_i$ is the natural-language specification (including I/O format and constraints) and $T_i$ is a deterministic test harness. Each problem is graded by executing the model-produced program against $T_i$, producing a binary correctness indicator. These tasks are instantiated as LeetCode problems with platform-provided acceptance criteria, and all generations are evaluated as Python solutions against Python-based submission templates.

\paragraph{Problem selection.} To reduce the risk of benchmark contamination from model training data, we select the $N=100$ most recently published problems available on the platform at collection time (newest first), excluding paid-only problems; no further curation is applied, so the benchmark is determined entirely by recency rank. This substantially reduces the probability of direct memorization relative to historically popular problems, but does not guarantee a contamination-free evaluation---recently published problems are often discussed online within days of release, and we do not perform contamination-detection analysis \citep{deng2023investigating}. We return to this in the Limitations section.

The dataset spans three difficulty tiers: 20 Easy, 50 Medium, and 30 Hard problems, covering a broad range of algorithmic topics including Array, Dynamic Programming, String, Graph, and Simulation. All problems satisfy the following properties:
\begin{itemize}
  \item \textbf{Deterministic grading:} each $T_i$ yields a deterministic accept/reject outcome.
  \item \textbf{Self-contained evaluation:} no network access or external dependencies are required at runtime.
  \item \textbf{Fixed snapshot:} all models are evaluated on the identical problem set to ensure comparability across runs.
  \item \textbf{No personal data:} problems consist of algorithmic specifications only and contain no personally identifying information.
\end{itemize}

In the evaluation framework, this benchmark is stored as a fixed local dataset snapshot rather than fetched online at run time. This design ensures that all model families are evaluated against the same problem set and ordering, and reduces the chance that later platform changes silently alter the benchmark during the experimental campaign.

\subsection{Models}

We evaluate 16 models across five provider families (OpenAI, Anthropic, Google, Qwen, DeepSeek), covering a range of model scales, training objectives, and architectural choices, including both standard and reasoning-specialized variants. Four are reasoning-specialized (o4-mini, Gemini 2.5 Flash+Think, QwQ-Plus, DeepSeek-R1). The full roster with types is given in Table~\ref{tab:models}, and exact API identifiers and access dates in Appendix~\ref{sec:model-snapshots}.

\subsection{Evaluation Framework}
For a model $m$ and a fixed evaluation configuration $c$ (prompt template, decoding parameters, toolchain), we run each problem $x_i$ multiple times. Each run produces a candidate program which is executed against $T_i$.

Let $Y_{i,r}^{(m,c)} \in \{0,1\}$ denote whether run $r \in \{1,\dots,R\}$ passes all tests for problem $i$. We treat compilation errors, runtime errors, timeouts, and wrong answers as failures ($0$). The pipeline separates generation from execution---outputs are generated via provider APIs, cached, normalized to a shared format, and only then submitted to the deterministic judge---which makes runs resumable and avoids conflating transient provider failures with execution outcomes.

\subsection{Repeated-Run Methodology}
We adopt a repeated-run design to estimate both mean performance and variability. For each $(m,c)$ pair we:
\begin{enumerate}
  \item Sample $R$ independent generations per problem under configuration $c$.
  \item Evaluate each generation deterministically against $T_i$.
  \item Aggregate outcomes into run-level and problem-level metrics with uncertainty estimates.
\end{enumerate}

The primary controlled intervention is the prompt template. We compare two prompt variants while holding the dataset, model, decoding parameters, and execution environment fixed. All models are evaluated with temperature $T=0.3$, top-$p=0.9$, and $R=5$ repeated runs per problem per prompt configuration. The default maximum output budget is $4{,}096$ tokens; reasoning-specialized models (o4-mini, Gemini 2.5 Flash+Think) use larger budgets to accommodate extended chain-of-thought output (see Appendix~\ref{sec:model-snapshots} for per-model details). The o-series API (o4-mini) does not support explicit temperature or top-$p$ settings; decoding parameters are omitted for that model (implications discussed in the Limitations section). We use a non-zero temperature because the goal of the study is to characterize repeated-run variability under realistic stochastic decoding rather than to collapse outputs toward near-deterministic behavior at $T=0$. The two prompt templates---\textsc{Detailed} and \textsc{Minimal}---are held fixed across all models to enable a controlled comparison of prompt sensitivity.

\paragraph{Prompt templates.} Both prompts instruct the model to return plain Python source code only and to follow the provided LeetCode method signature exactly. The \textsc{Detailed} prompt uses a longer structured template emphasizing careful constraint reading, edge-case coverage, submission readiness, and accepted-style performance. The \textsc{Minimal} prompt conveys the same task more compactly, with lighter instruction scaffolding. Because both templates share the same output constraints and code template, the prompt comparison is a controlled low-contrast comparison that isolates instruction density rather than changing the task itself. Accordingly, the study does not test broad prompt-engineering strategies; it evaluates a narrow prompt variation under otherwise fixed conditions.

\subsection{Metrics Definition}
\label{sec:metrics}

\paragraph{Run-Level Pass Rate (RLPR).}
$$
\mathrm{RLPR}(m,c)=\frac{1}{NR}\sum_{i=1}^{N}\sum_{r=1}^{R} Y_{i,r}^{(m,c)}.
$$
RLPR measures the probability that a \emph{random invocation} succeeds.

\paragraph{Perfect Stability Rate (PSR).}
$$
\mathrm{PSR}(m,c)=\frac{1}{N}\sum_{i=1}^{N} \mathbf{1}\!\left[\sum_{r=1}^{R} Y_{i,r}^{(m,c)}=R\right].
$$
PSR quantifies the fraction of tasks for which the system is reliably correct without retries.

\paragraph{Average Variance (AV).}
$\mathrm{AV}(m,c)=\frac{1}{N}\sum_{i=1}^{N} \frac{R}{R-1}\hat{p}_i(1-\hat{p}_i)$ with $\hat{p}_i=\frac{1}{R}\sum_r Y_{i,r}$ captures average instability across tasks.

\paragraph{Inference.}
We report 95\% Wilson score intervals for RLPR and PSR, and nonparametric bootstrap CIs for AV. Since both prompt templates are evaluated on the same problems, we compare PSR across prompt conditions using exact McNemar tests on the binary indicator of whether all $R=5$ runs pass. Reported statistics are descriptive summaries of the observed protocol rather than universal claims; prompt comparisons are specific to the two templates used here.

\section{Results}

\begin{table*}[t]
\centering
\small
\caption{Main results across all 16 models and two prompt configurations (100 problems, $R=5$ runs, 16{,}000 total evaluation instances). RLPR and PSR are in \%; AV is shown $\times 10^{-2}$ (i.e., multiply by 0.01 for the raw value). All entries are point estimates; 95\% CIs (Wilson score for RLPR/PSR; bootstrap for AV) available from authors on request. $\dagger$ reasoning-specialized model. $*$ output budget differs from the 4{,}096-token default: Gemini 2.5 Flash+Think used 24{,}576 tokens (8{,}192 thinking); o4-mini used 16{,}384 tokens. $\ddagger$ o4-mini decoding is not user-controllable; its PSR/AV may reflect reduced stochasticity rather than task reliability---see the Limitations section.}
\label{tab:overall-performance}
\setlength{\tabcolsep}{2.5pt}
\resizebox{\textwidth}{!}{
\begin{tabular}{llcccccc}
\toprule
\multirow{2}{*}{Family} & \multirow{2}{*}{Model}
  & \multicolumn{2}{c}{RLPR (\%)}
  & \multicolumn{2}{c}{PSR (\%)}
  & \multicolumn{2}{c}{AV ($\times 10^{-2}$)} \\
\cmidrule(lr){3-4}\cmidrule(lr){5-6}\cmidrule(lr){7-8}
& & D & M & D & M & D & M \\
\midrule
\multirow{3}{*}{OpenAI}
  & GPT-4.1               & 59.6 & 59.2 & 47.5 & 49.0 & 4.36 & 4.08 \\
  & GPT-4.1-mini          & 58.4 & 59.2 & 50.0 & 47.0 & 3.28 & 4.48 \\
  & o4-mini$^{\dagger\ddagger}$ & 86.2 & 85.6 & 75.0 & 69.0 & 3.36 & 5.20 \\
\midrule
\multirow{2}{*}{Anthropic}
  & Claude Sonnet 4.5 & 52.8 & 51.4 & 43.0 & 47.0 & 3.36 & 1.68 \\
  & Claude Haiku 4.5  & 41.2 & 40.8 & 35.0 & 34.0 & 2.40 & 2.00 \\
\midrule
\multirow{5}{*}{Google}
  & Gemini 3.1 Pro (preview)        & \textbf{86.2} & 85.8 & \textbf{75.0} & 74.0 & 3.36 & 3.92 \\
  & Gemini 3 Flash (preview)        & 80.8 & \textbf{83.2} & 67.0 & 70.0 & 5.28 & 4.40 \\
  & Gemini 3.1 Flash-Lite (preview) & 69.2 & 63.6 & 53.0 & 50.0 & 5.28 & 5.44 \\
  & Gemini 2.5 Flash$^*$            & 53.8 & 48.2 & 36.0 & 33.0 & 7.04 & 5.84 \\
  & Gemini 2.5 Flash+Think$^{\dagger*}$ & 71.6 & 74.4 & 58.0 & 57.0 & 4.64 & 4.72 \\
\midrule
\multirow{4}{*}{Qwen}
  & QwQ-Plus   & 72.8 & 72.6 & 59.0 & 58.0 & 4.80 & 4.72 \\
  & Qwen-Plus  & 63.4 & 61.2 & 52.0 & 51.0 & 4.72 & 4.64 \\
  & Qwen-Max   & 30.6 & 30.4 & 26.0 & 25.0 & 2.24 & 2.24 \\
  & Qwen-Turbo & 34.4 & 30.8 & 32.0 & 28.0 & 0.88 & 1.04 \\
\midrule
\multirow{2}{*}{DeepSeek}
  & DeepSeek-R1   & 84.2 & \textbf{84.4} & 74.0 & \textbf{75.0} & 3.76 & 3.20 \\
  & DeepSeek-V3.2 & 48.4 & 47.6 & 36.0 & 38.0 & 5.52 & 4.00 \\
\bottomrule
\end{tabular}
}
\end{table*}

\subsection{Overall Accuracy and Stability}
Table~\ref{tab:overall-performance} reports RLPR, PSR, and AV for all 32 model/prompt configurations. The top performers under RLPR are o4-mini (86.2\% D / 85.6\% M), Gemini 3.1 Pro (86.2\% D / 85.8\% M), and DeepSeek-R1 (84.2\% D / 84.4\% M). More broadly, the table shows that high run-level performance and high retry-free coverage often co-occur, but not in a one-to-one way: models with similar RLPR can still differ materially in PSR and AV.

Figure~\ref{fig:acc-stability-scatter} plots RLPR against PSR for all 32 configurations. All points fall below the diagonal, confirming that RLPR consistently overstates production reliability. Across all configurations, RLPR and PSR are strongly correlated ($r=0.985$, $p<0.001$), which is directionally expected because both metrics are derived from the same underlying per-problem success structure. The informative result is therefore not the existence of a high correlation itself, but the persistent and non-uniform gap between run-level success and retry-free coverage.

The gap $\Delta = \mathrm{RLPR} - \mathrm{PSR}$ is non-uniform and reaches its maximum in the mid-accuracy tier. Gemini 2.5 Flash has the largest gap (17.8\,pp under \textsc{Detailed}), followed by Gemini 3.1 Flash-Lite (16.2\,pp D) and Gemini 2.5 Flash+Think (17.4\,pp M). At the low end, Qwen-Turbo has a gap of only 2.4\,pp---not because it is uniformly reliable, but because its outcomes are highly polarized: many problems are either always failed or always passed, leaving less room for intermediate trial-to-trial variability. This pattern is reflected in AV: Qwen-Turbo has the lowest AV ($0.0088$), while Gemini 2.5 Flash has the highest ($0.0704$). Together, PSR and AV help distinguish between systems that are consistently strong, consistently weak, and unstable in the middle.

While RLPR $\geq$ PSR is mathematically expected for any non-trivial Bernoulli process, and Bernoulli variance peaks at $p \approx 0.5$ so mid-accuracy models exhibit the largest gaps by construction, the \emph{magnitude} of $\Delta$ and \emph{which} models fall in the mid-accuracy danger zone cannot be derived a priori, since per-problem success probabilities are unknown in advance. The empirical contribution is calibrating this expected pattern at scale: the gap is observable, model-specific, and meaningfully large (up to 17.8\,pp), so RLPR rankings alone cannot reveal how much of a model's accuracy translates into retry-free reliability.

\begin{figure}[t]
\centering
\includegraphics[width=0.92\linewidth]{figures/paper_figure1_accuracy_vs_stability.png}
\caption{RLPR vs.\ PSR scatter plot across all 32 configurations (16 models $\times$ 2 prompt types). Points are colored by model family and shaped by prompt type. The diagonal marks $\mathrm{PSR}=\mathrm{RLPR}$; points below indicate configurations where mean success overstates stable problem coverage.}
\label{fig:acc-stability-scatter}
\end{figure}

\subsection{Problem-Level Stability Analysis}
Aggregate rates hide qualitatively different reliability profiles. A cross-model stability heatmap (Figure~\ref{fig:stability-heatmap}) and the per-problem pass-rate distribution (Figure~\ref{fig:success-distribution}), both in Appendix~\ref{sec:extra}, show that models with similar mean pass rates can differ sharply in how their successes are distributed: a model with many $\hat{p}_i \approx 1$ and $\hat{p}_i \approx 0$ problems behaves very differently from one with many $\hat{p}_i \approx 0.4$--$0.6$ problems. Qwen-Turbo shows the most polarized profile (AV $=0.0088$), while Gemini 2.5 Flash has the highest AV ($0.0704$). Retry-based recovery does not eliminate this instability: PSR measures the fraction of tasks needing no retries, while AV quantifies how much of the remaining workload sits in a stochastic regime.

\subsection{Difficulty Gradient and Failure Modes}
\label{sec:difficulty}

Pooling outcomes by benchmark difficulty reveals a pronounced gradient: run-level pass rates are highest on Easy problems (90.6\%), lower on Medium (62.2\%), and substantially lower on Hard (34.0\%). The dominant non-accept verdict is \emph{Wrong Answer} ($\sim$58\%), followed by \emph{Time Limit Exceeded} ($\sim$27\%) and \emph{Runtime Error} ($\sim$11\%), with the remaining $\sim$4\% comprising compile errors and other verdicts. Most failures thus arise from incorrect algorithms rather than formatting noise, though timeouts and runtime failures increase on harder problems, indicating that instability is also tied to implementation robustness. A low PSR can therefore reflect two regimes---algorithmic inconsistency and implementation fragility---which imply different mitigation strategies.

\subsection{Reasoning vs.\ Standard Models (Exploratory)}
We tested whether reasoning-specialized models (DeepSeek-R1, QwQ-Plus, o4-mini, Gemini 2.5 Flash+Think) exhibit smaller accuracy--stability gaps after controlling for accuracy, by regressing $\Delta$ on RLPR and comparing group residuals. They do not: the mean residual for the eight reasoning configurations ($-0.04$) is statistically indistinguishable from the 24 standard configurations ($+0.01$; permutation test $p=0.97$), and mean AV is likewise comparable ($4.30$ vs.\ $3.81 \times 10^{-2}$). The pattern is model-specific---DeepSeek-R1 is markedly more stable than expected for its accuracy while Gemini 2.5 Flash+Think is less so---and severely underpowered at four models, so we read $p=0.97$ as ``insufficient evidence,'' not ``no effect.'' Full analysis and Figure~\ref{fig:reasoning-scatter} appear in Appendix~\ref{sec:extra}.

\subsection{Prompt Sensitivity (Exploratory)}
Comparing \textsc{Detailed} vs.\ \textsc{Minimal} prompting, effects are model-dependent and small: 11 of 16 models favor \textsc{Detailed} and 5 favor \textsc{Minimal}, with a mean absolute PSR difference of $\sim$2.3\,pp and no model reaching $p<0.05$ under exact McNemar tests \citep{zhuo2024prosa}. Because both templates share identical structure and differ only in instruction density, this is a deliberately low-contrast comparison and underpowered at $N=100$; meaningful effects cannot be ruled out. Full analysis and Figure~\ref{fig:prompt-comparison} appear in Appendix~\ref{sec:extra}.

\subsection{Accuracy--Stability Gap}
\label{sec:gap}

Beyond its magnitude, the gap $\Delta$ can reorder model rankings. Because $\Delta$ peaks in the mid-accuracy tier (as shown earlier), two models with similar RLPR can differ substantially in retry-free reliability, so a single RLPR ranking cannot infer the accuracy--stability gap; PSR makes the distinction explicit.

RLPR and PSR can disagree on model rankings despite the high correlation. For example, GPT-4.1 leads GPT-4.1-mini on RLPR (59.6\% vs.\ 58.4\%) but trails on PSR (47.5\% vs.\ 50.0\%)---a reversal driven by GPT-4.1's larger accuracy--stability gap (12.1\,pp vs.\ 8.4\,pp). This confirms that PSR captures deployment-relevant distinctions invisible to RLPR rankings alone.

\section{Conclusion}

We presented a repeated-run evaluation framework for deterministic programming tasks and used it to study the relationship between correctness and stability. Across 16 models, 100 problems, and 5 repeated runs per configuration (16{,}000 total instances), run-level pass rate and retry-free coverage are strongly correlated, but run-level accuracy consistently exceeds perfect stability. The accuracy--stability gap is non-uniform and reaches up to 17.8 percentage points, peaking in the mid-accuracy tier. Importantly, PSR and RLPR can reverse model rankings even at high correlation ($r=0.985$): GPT-4.1 leads GPT-4.1-mini on RLPR but trails on PSR, demonstrating that the two capture deployment-relevant distinctions invisible to single-run reporting. We view these metrics as complementary: RLPR characterizes average invocation success, while PSR and AV reveal whether that success is concentrated in consistently solved problems or dispersed across unstable outcomes. Prompt effects and the reasoning-vs-standard distinction are both heterogeneous and exploratory at the current sample size, suggesting that neither is a reliable lever for improving stability without model-specific tuning.

Future work should replicate the repeated-run analysis on non-coding deterministic NLP tasks (e.g., text-to-SQL) and control output token budgets to cleanly isolate the effect of reasoning on stability; broader orchestration and runtime dependencies require separate analysis \citep{jiang2026agenticaicybersecurityattack}. Taken together, these findings suggest that reporting the test--retest reliability of a score---not just its point value---is a necessary part of valid evaluation for any deterministic text-conditioned generation task, and a prerequisite for extending measurement to the harder non-stationary settings.

\section*{Limitations}

Several limitations qualify our results. First, the benchmark is small ($N=100$) and restricted to recent LeetCode-style problems: this reduces overlap with heavily circulated historical benchmarks but does not eliminate contamination risk, since recent contest problems are often discussed online within days of release. The size is also a tradeoff---16 models $\times$ 2 prompts $\times$ 5 runs still yields 16{,}000 judged instances, ample for broad correctness--stability patterns but underpowered for small prompt effects (consistent with the non-significant comparisons in Section~4.5). Second, the protocol is intentionally narrow: we vary prompt template and repeated sampling while holding the harness fixed, so it does not characterize broader prompt engineering, tool use, self-repair, retrieval, test-time selection, or multi-turn interaction. Third, access is mediated through provider APIs, so despite deterministic grading we cannot control backend changes or model refreshes; conclusions should be read relative to the documented snapshots and access dates, not as timeless statements about a provider. Fourth, our metrics compress distinct failure modes into a binary outcome and do not by themselves explain whether instability arises from reasoning errors, formatting mistakes, runtime failures, or judge edge cases. Fifth, a caveat applies to o4-mini: the o-series API exposes no temperature/top-$p$ controls, so if its effective decoding is near-deterministic, its PSR/AV measure the near-absence of stochasticity rather than stability under stochastic sampling, and should be interpreted cautiously. Finally, results are limited to Python code generation on deterministic tasks with executable tests, and should not be read as a complete account of reliability for other languages, open-ended generation, or production systems with additional orchestration layers.

\section*{Ethics Statement}

This paper evaluates commercial and open-access LLM APIs on deterministic programming tasks. The study does not involve human subjects, personal data collection, or annotation labor. The main ethical considerations are instead tied to platform use and deployment interpretation. First, the evaluation relies on automated submission to an external online judge. Such use should be understood as research benchmarking under the platform's operational constraints, and any public release should avoid exposing credentials, session artifacts, or procedures that would enable abusive automation. Second, the evaluated models are accessed through third-party APIs whose outputs and availability may be governed by provider-specific terms of service and may change over time. Our reported results should therefore be interpreted as measurements of accessed systems under a documented evaluation protocol, not as immutable properties of any provider's model family.

More broadly, stronger performance on coding benchmarks can increase both beneficial and harmful capabilities. Better repeated-run reliability can support legitimate software engineering workflows, but it may also lower the barrier to generating code in settings where security review is weak. For this reason, we frame our contribution as an evaluation methodology rather than as an endorsement of fully autonomous code deployment.

\bibliographystyle{colm2026_conference}
\bibliography{references}

\appendix

\section{Additional Stability Analyses}
\label{sec:extra}

This appendix contains the detailed problem-level, reasoning-vs-standard, and prompt-sensitivity analyses summarized in the Results section, together with their figures.

\subsection{Problem-Level Stability Detail}

Figure~\ref{fig:stability-heatmap} presents a cross-model stability overview. Each row corresponds to one of the 16 models, ordered by RLPR under \textsc{Detailed} prompting (highest at top). Each column corresponds to one of the 100 problems, grouped by difficulty tier (Easy, Medium, Hard) and sorted within each tier by decreasing mean pass rate across all models. Cell color encodes the empirical pass probability for that model--problem pair, averaged over $R=5$ runs: green indicates a problem solved consistently on every run, red indicates consistent failure, and yellow indicates high trial-to-trial variability.

\begin{figure}[t]
\centering
\includegraphics[width=0.96\linewidth]{figures/paper_figure2_stability_heatmap.png}
\caption{Cross-model stability heatmap (\textsc{Detailed} prompt). Rows: 16 models sorted by RLPR descending. Columns: 100 problems grouped by difficulty tier with dashed separators, and ordered within each tier by decreasing mean pass rate. Cell color encodes empirical pass rate over $R=5$ runs (green $=1.0$, red $=0.0$, yellow $\approx 0.5$).}
\label{fig:stability-heatmap}
\end{figure}

Figure~\ref{fig:success-distribution} shows the distribution of per-problem pass rates aggregated across all models. This view exposes qualitatively different reliability profiles that similar aggregate pass rates can hide: a model with many $\hat{p}_i \approx 1$ and $\hat{p}_i \approx 0$ problems behaves differently from a model with many $\hat{p}_i \approx 0.4$ or $0.6$ problems, even when their average run-level pass rate is similar.

\begin{figure}[t]
\centering
\includegraphics[width=0.92\linewidth]{figures/paper_figure3_success_distribution.png}
\caption{Problem-level success distribution (\textsc{Detailed} prompt). The figure contains one panel per model (16 panels total), arranged in decreasing order of PSR. Each panel shows the fraction of that model's 100 problems falling in each empirical pass-probability bin $\hat{p}_i \in \{0/5, 1/5, \ldots, 5/5\}$, where $0/5$ means the model failed all 5 runs and $5/5$ means it passed all 5.}
\label{fig:success-distribution}
\end{figure}

\subsection{Reasoning vs.\ Standard Models: Full Analysis}

Reasoning-specialized models generate an internal chain-of-thought before producing a final answer, which could plausibly improve output consistency by structuring the solution process. We test whether this architectural distinction manifests as a systematically smaller accuracy--stability gap after controlling for overall accuracy level.

\paragraph{Grouping.} We treat four models as reasoning-specialized (marked $\dagger$ in Table~\ref{tab:models}): DeepSeek-R1, QwQ-Plus, o4-mini, and Gemini 2.5 Flash+Think. The remaining 12 systems are classified as standard models. All conclusions from this analysis remain exploratory given the modest group sizes.

\paragraph{RLPR-controlled gap analysis.} Because $\Delta = \mathrm{RLPR} - \mathrm{PSR}$ is itself correlated with RLPR (higher-accuracy models tend to have larger absolute gaps, as shown in Section~\ref{sec:gap}), a raw comparison of mean gaps between groups would be confounded by accuracy level. To address this, we fit a linear model $\Delta \sim \mathrm{RLPR}$ across all 32 configurations ($R^2=0.40$, $p<0.001$) and examine the residuals: a negative residual indicates a smaller gap than expected at the model's accuracy level, i.e., greater stability relative to accuracy.

Figure~\ref{fig:reasoning-scatter} shows the result. Reasoning models show strongly heterogeneous behavior: DeepSeek-R1 falls well below the regression line (mean residual $\approx -4.4$), suggesting markedly better stability than expected for its accuracy level; o4-mini is near the line (residual $\approx -0.5$); while QwQ-Plus (residual $\approx +1.7$) and Gemini 2.5 Flash+Think (residual $\approx +3.0$) both fall above it, showing larger-than-expected gaps. Across all eight reasoning-model configurations, the mean residual is $-0.04$, compared with $+0.01$ for the 24 standard-model configurations---a difference of 0.05 percentage points. A permutation test (10{,}000 iterations, two-sided) yields $p=0.97$, providing no evidence that the residual difference is systematic.

\begin{figure}[t]
\centering
\includegraphics[width=0.90\linewidth]{figures/paper_figure5_reasoning_vs_standard.png}
\caption{Accuracy--stability gap ($\Delta$) vs.\ RLPR for all 32 configurations. Reasoning models (diamonds) and standard models (circles) are shown against an OLS fit. Points below the line have smaller gaps than expected at their accuracy level.}
\label{fig:reasoning-scatter}
\end{figure}

\paragraph{Variance comparison.} Average variance (AV) provides a complementary view of trial-to-trial variability. Across the eight reasoning-model configurations, the mean AV is $4.30 \times 10^{-2}$, compared with $3.81 \times 10^{-2}$ for the 24 standard-model configurations---again showing no systematic advantage for reasoning-specialized architectures on this metric.

\paragraph{Interpretation.} The four reasoning models do not exhibit systematically smaller accuracy--stability gaps or lower AV once accuracy is accounted for. DeepSeek-R1 shows suggestive evidence of above-expected stability (residual $\approx -4.4$\,pp), warranting verification in a larger sample. However, with only four reasoning models, the analysis is severely underpowered: the permutation test ($p=0.97$) should be read as ``insufficient evidence from a small sample'' rather than ``the effect is negligible.'' The pattern is model-specific: chain-of-thought reasoning neither reliably improves nor worsens stability within this sample. One important caveat: the Gemini 2.5 Flash vs.\ Flash+Think comparison is confounded by output token budget (4{,}096 vs.\ 24{,}576 tokens), so the pair should not be interpreted as a clean ablation of thinking mode alone.

\subsection{Prompt Sensitivity: Full Analysis}

Figure~\ref{fig:prompt-comparison} visualizes the prompt effect as a paired scatter plot: each of the 16 models appears as a single point with its \textsc{Minimal}-prompt PSR on the $x$-axis and \textsc{Detailed}-prompt PSR on the $y$-axis. Points above the diagonal indicate that detailed prompting improves stability.

\begin{figure}[t]
\centering
\includegraphics[width=0.8\linewidth]{figures/paper_figure4_prompt_comparison.png}
\caption{Prompt sensitivity: paired scatter of PSR under \textsc{Detailed} vs.\ \textsc{Minimal} prompting for all 16 models. Points above the diagonal indicate that detailed prompting improves stability. Points are colored by model family.}
\label{fig:prompt-comparison}
\end{figure}

Prompt effects are model-dependent: most models show small differences (typically within 1--3 percentage points of PSR), but Claude Sonnet 4.5 shows a reversal (higher PSR under \textsc{Minimal}), while Qwen-Turbo shows the opposite (4-point gain under \textsc{Detailed}). No single template uniformly dominates---consistent with prior findings on prompt sensitivity \citep{zhuo2024prosa}. Exact McNemar tests on the paired PSR indicators find no model reaching $p<0.05$; the largest PSR shift is 6 points (o4-mini, whose decoding regime is not user-controllable; see the Limitations section), and among the remaining 15 models the largest shift is 4 points ($p=0.219$). Across 16 models, 11 favor \textsc{Detailed} and 5 favor \textsc{Minimal} (mean absolute difference $\sim$2.3\,pp). This is a deliberately low-contrast comparison---both templates share identical structure, differing only in instruction density---so the null result may partly reflect insufficient contrast rather than a true absence of prompt effects. With $N=100$, a 4-point shift has less than 20\% power, meaning meaningful effects cannot be ruled out at the current sample size.

\section{Prompt Templates}

\subsection{Detailed Prompt}

\begin{lstlisting}
<role>
You are an expert Python algorithm engineer solving deterministic
LeetCode-style problems.
</role>

<task>
Produce a correct Python solution that fits the provided template and
is optimized for the problem constraints.
</task>

<instructions>
1. Read the full problem statement, constraints, and template before
   writing code.
2. Produce the best practical algorithm for the stated constraints,
   aiming for the optimal time complexity when possible.
3. Do not return a brute-force, quadratic, exponential, or placeholder
   solution when the constraints require a more efficient algorithm.
4. Use the required method signature from the template exactly.
5. Ensure the algorithm is complete and submission-ready, not a sketch
   or partial attempt.
6. Handle edge cases implied by the statement and constraints.
7. Return one complete Python solution that can be submitted directly.
</instructions>

<format_requirements>
Return plain Python source code only.
Start directly with the code.
Do not include explanations, markdown fences, or surrounding commentary.
The code must be valid, runnable Python 3 with no syntax errors.
The code must fit the provided code template without requiring manual edits.
The final code should be written for Accepted-style performance, not just
sample-case correctness.
Only include comments when they clarify non-obvious logic.
</format_requirements>

<problem>
{problem_description}
</problem>

<code_template>
{code_template}
</code_template>
\end{lstlisting}

\subsection{Minimal Prompt}

\begin{lstlisting}
<role>
You are a Python coding assistant solving a deterministic
LeetCode-style problem.
</role>

<task>
Write one correct and efficient Python solution that fits the provided
template.
</task>

<instructions>
1. Infer an algorithm that matches the constraints before writing code.
2. Prefer the optimal or near-optimal solution for the input limits.
3. Do not output brute-force code unless the constraints clearly make
   it acceptable.
4. Use the exact method signature from the template.
5. Return only final submission code.
</instructions>

<format_requirements>
Return plain Python source code only.
Start directly with the code.
Do not include explanations, markdown fences, or extra text.
The code must be valid, runnable Python 3 with no syntax errors.
The code must fit the provided code template without requiring manual edits.
The solution must be efficient enough to avoid obvious Time Limit
Exceeded outcomes under the given constraints.
</format_requirements>

<problem>
{problem_description}
</problem>

<code_template>
{code_template}
</code_template>
\end{lstlisting}

\section{Model Snapshots}
\label{sec:model-snapshots}

\begin{table}[t]
\centering
\small
\caption{Models evaluated in this study, grouped by provider family. $\dagger$ denotes reasoning-specialized models. Display names are editorial labels derived from API identifier strings; they are not official product version names. Exact API identifiers and access dates are listed in Table~\ref{tab:model-snapshots}. Google models marked ``(preview)'' were accessed via preview-tier API endpoints.}
\label{tab:models}
\resizebox{\linewidth}{!}{
\begin{tabular}{lll}
\toprule
Family & Model & Type \\
\midrule
\multirow{3}{*}{OpenAI}
  & GPT-4.1         & flagship \\
  & GPT-4.1-mini    & lightweight \\
  & o4-mini         & reasoning$^\dagger$ \\
\midrule
\multirow{2}{*}{Anthropic}
  & Claude Sonnet 4.5 & flagship \\
  & Claude Haiku 4.5  & lightweight \\
\midrule
\multirow{5}{*}{Google}
  & Gemini 3.1 Pro (preview)        & flagship \\
  & Gemini 3 Flash (preview)        & efficient \\
  & Gemini 3.1 Flash-Lite (preview) & lightweight \\
  & Gemini 2.5 Flash                & efficient \\
  & Gemini 2.5 Flash+Think          & reasoning$^\dagger$ \\
\midrule
\multirow{4}{*}{Qwen}
  & QwQ-Plus   & reasoning$^\dagger$ \\
  & Qwen-Plus  & flagship \\
  & Qwen-Max   & large \\
  & Qwen-Turbo & lightweight \\
\midrule
\multirow{2}{*}{DeepSeek}
  & DeepSeek-R1   & reasoning$^\dagger$ \\
  & DeepSeek-V3.2 & flagship \\
\bottomrule
\end{tabular}
}
\end{table}

\begin{table*}[t]
\centering
\small
\caption{Model snapshots used in the experiments. The API identifier column records the concrete model string used by the evaluation pipeline. Access dates are derived from the archived run dates for the corresponding completed paper runs.}
\label{tab:model-snapshots}
\resizebox{\textwidth}{!}{
\begin{tabular}{llll}
\toprule
Display Name & Provider & API Identifier & Access Date \\
\midrule
GPT-4.1 & OpenAI & gpt-4.1 & 2026-04-01 \\
GPT-4.1-mini & OpenAI & gpt-4.1-mini & 2026-03-31 \\
o4-mini & OpenAI & o4-mini & 2026-04-06 \\
Claude Sonnet 4.5 & Anthropic & claude-sonnet-4-5-20250929 & 2026-04-03 \\
Claude Haiku 4.5 & Anthropic & claude-haiku-4-5-20251001 & 2026-04-03 \\
Gemini 3.1 Pro (preview) & Google & gemini-3.1-pro-preview & 2026-03-28 \\
Gemini 3 Flash (preview) & Google & gemini-3-flash-preview & 2026-03-29 \\
Gemini 3.1 Flash-Lite (preview) & Google & gemini-3.1-flash-lite-preview & 2026-03-30 \\
Gemini 2.5 Flash & Google & gemini-2.5-flash & 2026-04-05 \\
Gemini 2.5 Flash+Think & Google & gemini-2.5-flash (thinkingBudget=8192) & 2026-04-05 \\
QwQ-Plus & Qwen & qwq-plus & 2026-03-23 \\
Qwen-Plus & Qwen & qwen-plus & 2026-03-29 \\
Qwen-Max & Qwen & qwen-max & 2026-03-23 \\
Qwen-Turbo & Qwen & qwen-turbo & 2026-03-23 \\
DeepSeek-R1 & DeepSeek & deepseek-r1 & 2026-03-24 \\
DeepSeek-V3.2 & DeepSeek & deepseek-v3.2 & 2026-03-25 \\
\bottomrule
\end{tabular}
}
\end{table*}

\end{document}
